% --- ARCHIVO: Anexos_Conceptos/anexo_cap3.tex ---

\chapter{Conceptos técnicos del Capítulo 3}
\label{anexo:conceptos_cap3}

\section{Impedancia de transferencia}
\label{anexo:impedancia}
En sistemas de potencia, la impedancia de transferencia entre dos nudos representa la oposición eléctrica al flujo de potencia entre ellos. Una alta impedancia de transferencia implica una red débilmente acoplada, en la que pequeñas variaciones de potencia inyectada pueden producir grandes variaciones de tensión y de ángulo de fase, deteriorando la firmeza del sistema.

\section{Oscilaciones forzadas y naturales}
\label{anexo:oscilaciones}
Una oscilación es \textit{forzada} cuando es inducida por una perturbación externa periódica ---típicamente un fallo o un comportamiento anómalo en el lazo de control de un equipo concreto---, frente a las \textit{oscilaciones naturales} o modos propios del sistema, cuya frecuencia viene determinada por la propia inercia y por las constantes electromecánicas de las máquinas síncronas conectadas.

\section{Ratio de amortiguamiento}
\label{anexo:amortiguamiento}
El \textit{ratio de amortiguamiento} (o amortiguamiento relativo) es un indicador adimensional que cuantifica la rapidez con la que una oscilación se atenúa tras una perturbación. Valores próximos al 5~\% se consideran un margen de seguridad operativo razonable en el sistema síncrono europeo; valores cercanos al 0~\% indican oscilaciones sostenidas, y valores negativos implican un crecimiento de la amplitud y, por tanto, un riesgo de inestabilidad.

\section{Efecto Ferranti}
\label{anexo:ferranti}
El \textit{Efecto Ferranti} describe el fenómeno por el cual, en una línea de transporte de alta tensión operada con poca o ninguna carga, la tensión en el extremo receptor supera a la del extremo emisor. La causa es la admitancia capacitiva distribuida de la línea: con flujo de potencia activa reducido, la carga capacitiva no se compensa con el consumo inductivo de las cargas, y el resultado es una sobretensión proporcional a la longitud de la línea. Es un fenómeno especialmente relevante al energizar líneas de 400~kV en vacío.

\section{Curvas de estabilidad de tensión Q-V}
\label{anexo:curvas_qv}
Las \textit{curvas Q-V} representan, para un nudo dado de la red, la relación entre la potencia reactiva inyectada o absorbida y la tensión resultante. La distancia entre el punto de operación y el punto de mínimo de la curva (\textit{nose point}) define el margen de estabilidad de tensión: cuanto menor sea ese margen, mayor será el riesgo de un colapso de tensión ante perturbaciones adicionales.

\section{Estabilizadores del Sistema de Potencia (PSS)}
\label{anexo:pss}
Los PSS (\textit{Power System Stabilizers}) son lazos de control adicionales instalados en el sistema de excitación de los grandes generadores síncronos que añaden amortiguamiento eléctrico a las oscilaciones electromecánicas del sistema. Cuando los grupos térmicos con PSS están desacoplados ---por ejemplo, durante un mediodía primaveral con baja demanda y alta producción solar---, la red pierde parte de la capacidad de amortiguamiento que estos dispositivos aportarían en condiciones de operación habitual.

\section{Cambiadores de Tomas en Carga (OLTC)}
\label{anexo:oltc}
Un Cambiador de Tomas en Carga (OLTC, \textit{On-Load Tap Changer}) es un mecanismo electromecánico instalado en los grandes transformadores de potencia que ajusta la relación de transformación ---y, por tanto, la tensión de salida del secundario--- sin interrumpir el flujo de energía. Regula la tensión ante variaciones lentas de carga, típicamente en un rango de $\pm 10~\%$ con escalones discretos. Su tiempo característico de respuesta, condicionado por la inercia mecánica del motor y los engranajes, es del orden de varios segundos por escalón.

\section{Sistema en por unidad (p.u.)}
\label{anexo:pu}
\textit{Per unit} (p.u.) es una notación adimensional empleada en sistemas de potencia que expresa cualquier magnitud (tensión, corriente, potencia, impedancia) como fracción de un valor de referencia (la base). Por convención, 1,0~p.u. equivale al valor nominal del equipo o nudo: una tensión de 1,2~p.u. en una red de 220~kV equivale, por tanto, a 264~kV.

\section{Bucle de retroalimentación (Feedback loop)}
\label{anexo:feedback_loop}
Un \textit{feedback loop} positivo, o bucle de retroalimentación positiva, describe un mecanismo en el que una perturbación inicial provoca una respuesta del sistema que, en lugar de oponerse al desvío, lo amplifica. En el contexto del incidente, cada disparo de planta IBR redujo la absorción de reactiva, lo que elevó la tensión, lo que a su vez provocó nuevos disparos: la respuesta del sistema reforzaba la perturbación en lugar de amortiguarla.

\section{Tasa de Cambio de Frecuencia (RoCoF)}
\label{anexo:rocof}
La RoCoF (\textit{Rate of Change of Frequency}) es la derivada temporal de la frecuencia del sistema, $df/dt$. Mide la velocidad con la que la frecuencia se desvía de su valor nominal (50~Hz) ante un desequilibrio de potencia activa. En sistemas con inercia elevada, la RoCoF es moderada; en sistemas con baja inercia, los mismos desequilibrios de potencia generan tasas de variación mucho mayores, comprometiendo la actuación coordinada de los sistemas de protección.